% This file is part of the Open Source project 'proTirone'
% (c) 2025 Karsten Reincke (https://github.com/protirone/protico.ltx)
% It is distributed under the terms of the creative commons license
% CC-BY-4.0 (= https://creativecommons.org/licenses/by/4.0/)

\selectlanguage{ngerman}
\begin{frame}  
  \frametitle{LF.CX:AP2-Strategie:Vorbereitung}

  \vspace{0.5cm}
  \begin{itemize}
    \item  einen knalligen Textmarker - für \textcolor{violet}{die eigentlichen Aufgaben}
  \end{itemize}
  \begin{center}
    
\includegraphics[height=0.8cm]{\imgLf/textmarker-violett.png}
  \end{center}

  \vspace{0.5cm}
  \begin{itemize}
    \item  einen positiv-dezenten Textmarker - für \textcolor{cyan}{direkte Ergänzungen}   
  \end{itemize}
  \begin{center}
    
\includegraphics[height=0.8cm]{\imgLf/textmarker-turquoise.png}
  \end{center}
 
  \vspace{0.5cm}
  \begin{itemize}
    \item  einen positiv-sanften Textmarker - für  \textcolor{teal}{Hintergrundinformationen}
  \end{itemize}
  \begin{center}
  
\includegraphics[height=0.8cm]{\imgLf/textmarker-green.png}
  \end{center} 
  
  \vspace{0.5cm}
  \begin{tiny}
  \begin{flushright}
    es gibt noch viele andere \href{https://www.amazon.de/s?k=pastelle+textmarker}{tolle Textmarker}
  \end{flushright}
  \end{tiny}

\end{frame}


\begin{frame}  
  \frametitle{LF.CX:AP2-Strategie:Prozess}

  \vspace{0.5cm}
  \begin{enumerate}
    \item \textbf{\textcolor{violet}{Grobscan}}: Jeden \textcolor{violet}{Aufgabensatz* knallig} einfärben.
    \item Feinklärung zu jedem Aufgabensatz:
    \begin{itemize}
      \item {Direkte Ergänzungen} dazu \textcolor{cyan}{positiv-dezent} einfärben.
      \item Entferntere \textcolor{teal}{Hintergrundinformationen} dazu \textcolor{teal}{positiv-sanft} markieren.
    \end{itemize}
    
    \item Quick-Wins einsammeln.
    \item Dicke Bretter bohren.
  \end{enumerate}

  
  \vspace{0.5cm}
  \begin{tiny}
  \begin{flushright}
    *) Aufgabensätze beginnen mit \textit{[Ergänzen | Erläutern | Beschreiben | Berechnen | Benennen | Ordnen | ... ] Sie ... }
  \end{flushright}
  \end{tiny}

  \vspace{2cm}
  \begin{tiny}
  \begin{flushright}
    Erprobt und verfeinert von/mit \textit{GSLDK 12ip/iv23}
  \end{flushright}
  \end{tiny}
\end{frame}

\begin{frame}  
  \frametitle{LF.CX:AP2-Strategie:Prozess:BPMN}
  \begin{center}
  \includegraphics[height=7.5cm]{\imgGl/apx-process.png}
  \end{center} 

\end{frame}

\begin{frame}  
  \frametitle{LF.CX:AP2-Strategie:Ergebnis}

  \begin{footnotesize}
    Die Kreincke-Hard-GmbH möchte \textcolor{teal}{betriebliche Abläufe der mobilen Erfassungssysteme anpassen} und beginnt dafür Projekt AWG ('Alles wird gut'). Der Scope des Projektes macht es nötigt, \enquote{sich mit verschiedenen Frage- und Problemstellungen bereits frühzeitig zu beschäftigen}.

  \begin{itemize}
    \item[a] Für die That's-it-Analyse sollen sie 'verschiedenen Eigenschafts- und \textcolor{cyan}{Unterscheidungsmerkmale der bei den Kunden im Umlauf befindlichen Geräte} erfassen'
    \begin{itemize}
      \item[aa] \textcolor{violet}{Beschreiben Sie 3 Verfahren, die eine That's-it-Beschreibung begünstigen}
      \item[...]
    \end{itemize}
    \item[...]
  \end{itemize}
  
  \end{footnotesize}

  
  \vspace{0.5cm}
  \begin{tiny}
  \begin{flushright}
    Die Zitate referieren lose auf AP2-FIDPA-Sommer-2024-GA1. Mögliche Antworten wären:\\ \textit{Interviews}, \textit{Fragebögen}, \textit{Beobachtungen}
  \end{flushright}
  \end{tiny}

\end{frame}